
\documentclass[a4paper]{article}
\usepackage[pdftex]{graphicx}
\usepackage[utf8]{inputenc}
\usepackage{enumerate}
\usepackage{siunitx}
\sisetup{locale=DE} 
\usepackage{href-ul}
\hypersetup{
	colorlinks=true,
	linkcolor=blue,
	urlcolor=blue}
\usepackage{icomma}
\usepackage{amssymb}
\usepackage{geometry}
\geometry{a4paper, top=15mm, left=15mm, right=15mm, bottom=15mm,
	headsep=10mm, footskip=12mm}
%Source: img_0235

% Start the document
\begin{document}
%Titel
{\bf Übungsblatt Schwingkreis }\\
\begin{enumerate}[1.]
%Aufgabe 1
\item Eine Spule mit der Induktivität $L$ und ein Kondensator mit der Kapazität $C$ bilden einen Schwingkreis. Der ohmsche Widerstand des Schwingkreises sei vernachlässigbar klein. Die Gesamtenergie des Schwingkreises beträgt $W_{ges}=\SI{2,9e-3}{\joule}$. Für die Stromstärke
$I$ im Schwingkreis zu einem Zeitpunkt $t$ mit $ t\ge 0~s$ gilt: 
$ I(t)= -\hat{I}\cdot \sin{ \left( {2 \pi \over T} \cdot t \right)}$, wobei $\hat{I}=\SI{120}{\milli\ampere}$ der Scheitelwert der Stromstärke und $T=\SI{2,0}{\milli\second}$ die Periodendauer der elektromagnetischen Schwingung ist.

\begin{enumerate}[a)]
\item Stellen Sie den zeitlichen Verlauf der Stromstärke $I$ für $0 \le t \le \SI{2,0}{\milli\second}$ in einem $t-I$-Diagramm dar.
\item Bei der elektromagnetischen Schwingung wird der Schwingkreiskondensator ständig aufgeladen und entladen. Geben Sie ein Zeitintervall an, in dem der Kondensator aufgeladen wird. Begründen Sie Ihre Antwort.
\item Berechnen Sie die Induktivität $L$ der Spule und die Kapazität $C$  des Kondensators. $[\textnormal{Teilergebnis:} \quad L=\SI{0,40}{\henry}]$
\item Kennzeichnen Sie im Diagramm von Teilaufgabe 1. a) den maximalen Betrag $Q_{max}$ der Ladung, die der Kondensator aufnimmt. Berechnen Sie  $Q_{max}$.
\item Berechnen Sie die im Kondensator gespeicherte elektrische Energie $W_{el,1}$ für einen Zeitpunkt $t_1$, zu dem die Stromstärke $I$ im Schwingkreis den Wert $I=\SI{90}{\milli\ampere}$ annimmt.
\end{enumerate}

%Aufgabe 2
\item Ein Schwingkreis regt einen Sendedipol $D_S$, dessen Länge $l$ verändert werden kann, zu Schwingungen mit der Frequenz $f=\SI{2,5}{\giga\hertz}$ an. Bei der Länge $l_0$ wird der Sendedipol in der Grundschwingung angeregt. 
\begin{enumerate}[a)]
\item Berechnen Sie $l_0$.

Der Sendedipol $D_S$ und ein baugleicher Empfangsdipol $D_E$ werden senk\-recht zu einer Horizontalebene angeordnet, wobei diese Horizontalebene durch die beiden Dipolmitten verläuft. Der Empfangsdipol registriert die Schwing\-ungsamplitude $E_0$ und empfängt die Leistung $P_0$. Zwischen $D_S$ und $D_E$ wird ein Hertz'sches Gitter so aufgestellt, dass die Gitterebene senkrecht zur Ge\-raden durch die beiden Dipolmitten ausgerichtet ist. Die Gitterstäbe sind wesentlich länger als die Dipole $D_S$ und $D_E$. Zunächst verlaufen die Gitterstäbe parallel zu den beiden Dipolen.

\item Beschreiben und Begründen Sie die Auswirkungen des Gitters auf die elektromagnetische Welle vor und hinter dem Gitter.

\item Das Gitter wird nun um den Winkel $\alpha=55^\circ$ gedreht, wobei die Gerade durch die beiden Dipolmitten die Drehachse ist. Berechnen Sie, wieviel Prozent der Leistung $P_0$ der Dipol $D_E$ jetzt empfängt. Veranschaulichen Sie Ihre Überlegungen anhand einer Skizze. 
\end{enumerate}

\end{enumerate} 

\fbox{
	\begin{minipage}{0.5\textwidth}
		Zur Lösung bitte \href{https://www.okuyakl.de/phys/p13schikL235/ll235.pdf}{hier klicken} oder den QR-Code scannen.\\
	Weitere Arbeitsblätter gibt es unter 
	
	\href{https://www.okuyakl.de}{www.okuyakl.de}
	\end{minipage}
	\hfill
	\begin{minipage}{0.4\textwidth}
		\includegraphics[width=1.5 cm]{../../viecher/zwe03}
		\includegraphics[width=3 cm]{qr235}
		\includegraphics[width=2 cm]{../../viecher/afanticon1}
		
	\end{minipage}}

\end{document}%Lösung----------------------------------------------
